% Generated by Listas/make_statements.py from pratica-lecture-08.md, item 24. Do not edit.
\textbf{Source.} Mas-Colell, Whinston \& Green (1995), \emph{Microeconomic Theory}, Ch. 13 Exercises, Exercise 13.C.2\^{}C, p.~475 {[}PDF p.~248{]}. Reproduced verbatim. Answer: none in the book; entries for 13.C.1, 13.C.2, 13.C.4 and 13.C.5 are in the MWG solutions manual (\texttt{Livros/}).



\begin{quote}

\textbf{13.C.2\^{}C} Reconsider the two-type signaling model with \(r(\theta_L) = r(\theta_H) = 0\), assuming a worker's productivity is \(\theta(1 + \mu e)\) with \(\mu > 0\). Identify the separating and pooling perfect Bayesian equilibria, and relate them to the perfect information competitive outcome.

\end{quote}



\textbf{Model of Section 13.C (pp.~450--451 {[}PDF p.~236{]}), copied from the book because the exercise refers to it.}



\begin{quote}

Consider the following adaptation of the model discussed in Section 13.B. For simplicity, we restrict attention to the case of two types of workers with productivities \(\theta_H\) and \(\theta_L\), where \(\theta_H > \theta_L > 0\) and \(\lambda = \text{Prob}\,(\theta = \theta_H) \in (0, 1)\). The important extension of our previous model is that before entering the job market a worker can get some education, and the amount of education that a worker receives is observable. To make matters particularly stark, we assume that education does \emph{nothing} for a worker's productivity (see Exercise 13.C.2 for the case of productive signaling). The cost of obtaining education level \(e\) for a type \(\theta\) worker (the cost may be of either monetary or psychic origin) is given by the twice continuously differentiable function \(c(e, \theta)\), with \(c(0, \theta) = 0\), \(c_e(e, \theta) > 0\), \(c_{ee}(e, \theta) > 0\), \(c_\theta(e, \theta) < 0\) for all \(e > 0\), and \(c_{e\theta}(e, \theta) < 0\) (subscripts denote partial derivatives). Thus, both the cost and the marginal cost of education are assumed to be lower for high-ability workers; for example, the work required to obtain a degree might be easier for a high-ability individual. Letting \(u(w, e \mid \theta)\) denote the utility of a type \(\theta\) worker who chooses education level \(e\) and receives wage \(w\), we take \(u(w, e \mid \theta)\) to equal her wage less any educational costs incurred: \(u(w, e \mid \theta) = w - c(e, \theta)\). As in Section 13.B, a worker of type \(\theta\) can earn \(r(\theta)\) by working at home.



A portion of the game tree for this model is shown in Figure 13.C.1. Initially, a random move of nature determines whether a worker is of high or low ability. Then, conditional on her type, the worker chooses how much education to obtain. After obtaining her chosen education level, the worker enters the job market. Conditional on the observed education level of the worker, two firms simultaneously make wage offers to her. Finally, the worker decides whether to work for a firm and, if so, which one.

\end{quote}

